\section*{الفصل \ref{ch:g12:ph-conductivity-titrations} --- المعايرة بقياس pH وبقياس الموصلية}

\begin{solution}{exo:g12:ph-conductivity-titrations:1}
$V_E = \qty{20.0}{mL}$؛ $c = 0.100 \times 20.0 / 20.0 = \qty{0.100}{mol/L}$.
\end{solution}

\begin{solution}{exo:g12:ph-conductivity-titrations:2}
عند \omterm{def:g12:ph-conductivity-titrations:ph-metric}{نصف التكافؤ} \omterm{def:g9:acids-bases-ph:ph}{pH} $= pK_a$: $pK_a = 3.75$، وهي قريبة من قيمة حمض الميثانويك
(3.74).
\end{solution}

\begin{solution}{exo:g12:ph-conductivity-titrations:3}
الفينول فثالين: فمجاله 8.0--10.0 يقع داخل القفزة، وهي قاعدية.
أما أحمر الميثيل (4.2--6.3) فيتغير لونه قرب \omterm{def:g12:ph-conductivity-titrations:ph-metric}{نصف التكافؤ}، مبكرًا جدًا.
\end{solution}

\begin{solution}{exo:g12:ph-conductivity-titrations:4}
تنحرف استجابة المسبار مع الزمن ودرجة الحرارة: فيُضبط ليقرأ
قراءة صحيحة في محلولين منظِّمين معروفي قيمة \omterm{def:g9:acids-bases-ph:ph}{pH}.
\end{solution}

\begin{solution}{exo:g12:ph-conductivity-titrations:5}
الترتيب تنازليًا: \ce{H3O+} (349.8)، ثم \ce{OH-} (198.3)، ثم \ce{Cl-} (76.2)، ثم \ce{Na+} (50.3).
\end{solution}

\begin{solution}{exo:g12:ph-conductivity-titrations:6}
$c = 0.050 \times 14.6 / 20.0 = \qty{0.0365}{mol/L}$.
\end{solution}

\begin{solution}{exo:g12:ph-conductivity-titrations:7}
الميول: 2.0، و4.5، و17.5 و\qty{2.5}{mL^{-1}}. وأكبرها بين
10.0 و\qty{10.2}{mL}: $V_E \approx \qty{10.1}{mL}$.
\end{solution}

\begin{solution}{exo:g12:ph-conductivity-titrations:8}
قبل \omterm{def:g11:titration:equivalence}{التكافؤ}، كل \omterm{def:g9:ions:ion}{أيون} \ce{OH-} يضاف يزيل \omterm{def:g9:ions:ion}{أيون} \ce{H3O+} (شديد
التوصيل) ويأتي \omterm{def:g9:ions:ion}{بأيون} \ce{Na+} (أقل توصيلًا بكثير): فتنخفض
\omterm{def:g12:ph-conductivity-titrations:conductivity}{الموصلية}. وبعده، تتراكم \omterm{def:g9:ions:ion}{أيونات} \ce{Na+} و\ce{OH-} دون أن
تتفاعل: فترتفع.
\end{solution}

\begin{solution}{exo:g12:ph-conductivity-titrations:9}
في البداية، قليل من \omterm{def:g9:ions:ion}{الأيونات} (\omterm{def:g12:ka-and-pka:strong-acid}{فالحمض الضعيف} لا يكاد يتفاعل مع الماء): \omterm{def:g12:ph-conductivity-titrations:conductivity}{موصلية}
منخفضة. وقبل \omterm{def:g11:titration:equivalence}{التكافؤ}، تتكوّن \omterm{def:g9:ions:ion}{أيونات} الإيثانوات والصوديوم
بدل \omterm{def:g7:atoms-and-molecules:molecule}{جزيئات} غير مشحونة: فترتفع، ببطء. وبعد \omterm{def:g11:titration:equivalence}{التكافؤ}، تتراكم
\ce{Na+} و\ce{OH-}: فترتفع أسرع. مستقيمان صاعدان،
والثاني أشد انحدارًا.
\end{solution}

\begin{solution}{exo:g12:ph-conductivity-titrations:10}
نحو 2.9 و4.76 و8.7. فحمض الإيثانويك ضعيف: لا يعطي \omterm{def:g9:ions:ion}{أيونات} أكسونيوم
إلا جزء ضئيل منه، فتبدأ قيمة \omterm{def:g9:acids-bases-ph:ph}{pH} أعلى من القيمة 1.0 لحمض الهيدروكلوريك
بتركيز مماثل.
\end{solution}

\begin{solution}{exo:g12:ph-conductivity-titrations:11}
لكي لا يكاد حجم \omterm{def:g11:titration:titration}{المحلول المعايِر} المضاف يغيّر الحجم الكلي: فتتغير
التراكيز، ومن ثم \omterm{def:g12:ph-conductivity-titrations:conductivity}{الموصلية}، على خطوط مستقيمة.
\end{solution}

\begin{solution}{exo:g12:ph-conductivity-titrations:12}
تتفاعل \omterm{def:g9:ions:ion}{أيونات} الهيدروكسيد أولًا مع \omterm{def:g12:ka-and-pka:strong-acid}{الحمض القوي} (\ce{H3O+})، ثم
مع الضعيف. فالقفزة الأولى تشير إلى نهاية حمض الهيدروكلوريك
(\omterm{def:g11:titration:equivalence}{وحجم التكافؤ} الأول يقيسه)، والثانية إلى نهاية حمض الإيثانويك
(والحجم بين القفزتين يقيسه).
\end{solution}

\begin{solution}{exo:g12:ph-conductivity-titrations:13}
البداية: $(349.8 + 76.2) \times 0.0100 = \qty{4.26}{mS/cm}$. وعند
\omterm{def:g11:titration:equivalence}{التكافؤ}: $[\ce{Na+}] = [\ce{Cl-}] = 1.00/110 = \qty{0.00909}{mol/L}$،
إذن $(50.3 + 76.2) \times 0.00909 = \qty{1.15}{mS/cm}$. وكلتاهما كما في
الشكل.
\end{solution}

\begin{solution}{exo:g12:ph-conductivity-titrations:14}
لا خطأ في $V_E$: فالقفزة تقع عند الحجم نفسه، مهما كان انزياح
القراءات. أما $pK_a$ المقروءة عند \omterm{def:g12:ph-conductivity-titrations:ph-metric}{نصف التكافؤ} فأعلى بمقدار 0.2.
\end{solution}

\begin{solution}{exo:g12:ph-conductivity-titrations:15}
لا تتغير قيمة \omterm{def:g9:acids-bases-ph:ph}{pH} خلال هذا الترسيب، لكن \omterm{def:g9:ions:ion}{الأيونات} تتغير:
فقبل \omterm{def:g11:titration:equivalence}{التكافؤ} كل \omterm{def:g9:ions:ion}{أيون} \ce{Ag+} يضاف يزيل \omterm{def:g9:ions:ion}{أيون} \ce{Cl-} ويأتي \omterm{def:g9:ions:ion}{بأيون}
\ce{NO3-} ذي \omterm{def:g12:ph-conductivity-titrations:conductivity}{موصلية} مماثلة، فتبقى \omterm{def:g12:ph-conductivity-titrations:conductivity}{الموصلية} شبه
ثابتة (\omterm{def:g9:ions:ion}{فأيونات} الصوديوم موجودة منذ البداية)؛ وبعد \omterm{def:g11:titration:equivalence}{التكافؤ} تتراكم \ce{Ag+} و\ce{NO3-}
فترتفع. خط أفقي، ثم خط صاعد.
\end{solution}

\begin{solution}{pb:g12:ph-conductivity-titrations:1}
\textbf{1.} \qty{6}{g} من حمض الإيثانويك في كل \qty{100}{g} من الخل.

\textbf{2.} يحوي \qty{101}{g} من الخل $6.0 \times 101/100 =
\qty{6.06}{g}$؛ و$M = \qty{60.0}{g/mol}$: أي \qty{0.101}{mol} في
\qty{0.1000}{L}، أي \qty{1.01}{mol/L}.

\textbf{3.} لو لم يُخفَّف، لتطلب \qty{10.0}{mL} منه نحو \qty{101}{mL} من
هيدروكسيد الصوديوم بتركيز \qty{0.100}{mol/L}: أكثر مما تسع السحاحة.

\textbf{4.} ماصة معايرة سعتها \qty{10.0}{mL} وحوجلة معايرة سعتها \qty{100.0}{mL}.

\textbf{5.} \ce{CH3COOH + OH- -> CH3COO- + H2O}.

\textbf{6.} $V_E = \qty{10.1}{mL}$.

\textbf{7.} $c = 0.100 \times 10.1 / 10.0 = \qty{0.101}{mol/L}$.

\textbf{8.} $10 \times 0.101 = \qty{1.01}{mol/L}$.

\textbf{9.} \omterm{def:g9:acids-bases-ph:ph}{pH} 4.76 عند \qty{5.05}{mL}: وهي قيمة $pK_a$ حمض الإيثانويك.

\textbf{10.} عند \omterm{def:g11:titration:equivalence}{التكافؤ} يحوي \omterm{def:g2:mixing-and-dissolving:dissolve}{المحلول} إيثانوات الصوديوم، والإيثانوات
\omterm{def:g12:ka-and-pka:strong-acid}{قاعدة ضعيفة}.

\textbf{11.} الفينول فثالين، الذي يحوي مجاله 8.0--10.0 قيمة \omterm{def:g9:acids-bases-ph:ph}{pH} عند
\omterm{def:g11:titration:equivalence}{التكافؤ} (نحو 8.7).

\textbf{12.} \omterm{def:g9:ions:ion}{أيونات} الصوديوم \omterm{def:g9:ions:ion}{وأيونات} الإيثانوات.

\textbf{13.} كل \omterm{def:g9:ions:ion}{أيون} \ce{OH-} يضاف يحوّل جزيئًا شبه غير متأين إلى
\ce{CH3COO-}، مع \omterm{def:g9:ions:ion}{أيون} \ce{Na+}: فتضاف \omterm{def:g9:ions:ion}{أيونات} متواضعة التوصيل
ببطء.

\textbf{14.} بعد \omterm{def:g11:titration:equivalence}{التكافؤ}، تتراكم \omterm{def:g9:ions:ion}{أيونات} الهيدروكسيد، الشديدة التوصيل،
مع \omterm{def:g9:ions:ion}{أيونات} الصوديوم.

\textbf{15.} \omterm{def:g12:ka-and-pka:strong-acid}{الحمض الضعيف} لا يعطي الماء إلا قليلًا جدًا من \omterm{def:g9:ions:ion}{الأيونات}.

\textbf{16.} عند $V_E = \qty{10.1}{mL}$، كما في \omterm{def:g12:ph-conductivity-titrations:ph-metric}{المعايرة بقياس pH}.

\textbf{17.} $1.01 \times 60.0 = \qty{60.6}{g/L}$.

\textbf{18.} \qty{6.06}{g} في \qty{100.0}{mL}.

\textbf{19.} $6.06 \times 100 / 101 = \qty{6.0}{g}$ في كل \qty{100}{g}.

\textbf{20.} حموضة الخل \qty{6.0}{\percent}: فاللصاقة
صحيحة.
\end{solution}
